\section*{अध्याय \ref{ch:g10:freshwater-resources} --- मीठे जल के संसाधन}

\begin{solution}{exo:g10:freshwater-resources:1}
\omterm{def:g10:freshwater-resources:stock}{भंडार} किसी दिए हुए समय पर किसी जलाशय में रखे जल की मात्रा है,
\unit{km^3} में (या किलोग्राम में)। \omterm{def:g10:freshwater-resources:stock}{अभिवाह} वह मात्रा है जो प्रति इकाई समय एक
जलाशय से दूसरे में जाती है, \unit{km^3} प्रति वर्ष में। \omterm{def:g10:freshwater-resources:residence-time}{निवास
काल} वह औसत समय है जो जल किसी जलाशय में बिताता है, दिनों या वर्षों में;
स्थायी अवस्था में यह \omterm{def:g10:freshwater-resources:stock}{भंडार} को \omterm{def:g10:freshwater-resources:stock}{अभिवाह} से भाग देने पर मिलता है।
\end{solution}

\begin{solution}{exo:g10:freshwater-resources:2}
सारे जल का लगभग \qty{2.5}{\percent} मीठा है। मीठे जल का
\qty{68.7}{\percent} हिम-टोपियों, हिमनदों और स्थायी हिम में जमा है, और
\qty{30.1}{\percent} भूमिगत है। नदियाँ सबसे कम रखती हैं,
\num{2120}~\unit{km^3}।
\end{solution}

\begin{solution}{exo:g10:freshwater-resources:3}
\omterm{def:g10:freshwater-resources:renewable}{नवीकरणीय मीठा जल} वह मीठा जल है जो हर वर्ष वर्षण द्वारा नवीकृत होता है
और किसी क्षेत्र से अपवाह तथा भौम जल के रूप में बहता है: प्रति वर्ष एक आयतन,
यानी \omterm{def:g10:freshwater-resources:stock}{अभिवाह}। \omterm{def:g10:freshwater-resources:fossil-groundwater}{जीवाश्म भौम जल} वह जल है जो बहुत पहले किसी जलभृत में प्रवेश कर चुका था और जिसे
आज का पुनर्भरण नवीकृत नहीं करता: यह ऐसा संचय है जिसे केवल घटाया जा सकता है,
यानी \omterm{def:g10:freshwater-resources:stock}{भंडार}।
\end{solution}

\begin{solution}{exo:g10:freshwater-resources:4}
$91\,000/2120 \approx 43$ गुना अधिक।
\end{solution}

\begin{solution}{exo:g10:freshwater-resources:5}
वायुमंडल अपना जल सबसे तेज़ी से नवीकृत करता है (लगभग \num{9.5} दिन), परिकलित जलाशयों में
महासागर सबसे धीरे (लगभग \num{3200} वर्ष); गहरा
भौम जल इससे भी पुराना है। $3200 \times 365.25/9.5 \approx \num{120000}$
गुना अधिक देर।
\end{solution}

\begin{solution}{exo:g10:freshwater-resources:6}
आगत: $413\,000 + 73\,000 = \num{486000}$~\unit{km^3} प्रति वर्ष, इसलिए
$T = 12\,700/486\,000 \approx \num{0.0261}$ वर्ष $\approx \num{9.5}$ दिन।
निर्गत: $373\,000 + 113\,000 = \num{486000}$~\unit{km^3}, उतना ही: वायु
स्थायी अवस्था में है, और \omterm{def:g10:freshwater-resources:residence-time}{निवास काल} वही है।
\end{solution}

\begin{solution}{exo:g10:freshwater-resources:7}
$1\,335\,040/413 \approx \num{3230}$ वर्ष। स्थायी अवस्था में महासागर
उतना ही पाता है जितना खोता है: $373 + 40 = 413$ हज़ार~\unit{km^3} प्रति वर्ष,
इसलिए आगतों से परिकलित \omterm{def:g10:freshwater-resources:residence-time}{निवास काल} वही है।
\end{solution}

\begin{solution}{exo:g10:freshwater-resources:8}
स्थल प्रति वर्ष \num{113000}~\unit{km^3} पाता है; \num{73000}
वाष्पन और वाष्पोत्सर्जन द्वारा वायु में लौटते हैं और \num{40000} बहकर समुद्र तक जाते हैं:
$73\,000 + 40\,000 = \num{113000}$। वायु में लौटने वाला अंश
$73/113 \approx \qty{65}{\percent}$ है।
\end{solution}

\begin{solution}{exo:g10:freshwater-resources:9}
$273.2 \times 10^6 \times 1233.5 \approx \qty{3.37e11}{m^3}$, यानी
\qty{337}{km^3}: लगभग उतना जितना सारा जल अरल सागर में
1987 में बचा था।
\end{solution}

\begin{solution}{exo:g10:freshwater-resources:10}
$15.8 \times 0.3048 \approx \qty{4.82}{m}$। 1950 से
2015 तक के 65 वर्षों में, $482/65 \approx \qty{7.4}{cm}$ प्रति वर्ष।
\end{solution}

\begin{solution}{exo:g10:freshwater-resources:11}
$2/0.01 = 200$ वर्ष। \omterm{def:g10:freshwater-resources:stock}{भंडार} अब प्रति वर्ष $0.05 - 0.01 =
\qty{0.04}{km^3}$ घटता है, इसलिए वह $2/0.04 = 50$ वर्ष टिकता है।
\end{solution}

\begin{solution}{exo:g10:freshwater-resources:12}
1960: $1093/71.1 \approx 15$ वर्ष; 2010: $98.1/11.4 \approx 8.6$ वर्ष।
वाष्पन झील के हर वर्ग किलोमीटर से हर वर्ष जल की लगभग उतनी ही मोटी परत
निकालता है; झील के सिकुड़ने के साथ उसका क्षेत्रफल सिकुड़ा, और उसके साथ
वाष्पित कुल आयतन भी।
\end{solution}

\begin{solution}{exo:g10:freshwater-resources:13}
(1) शुद्ध हानि $1.2 - 0.2 = \qty{1.0}{m}$ प्रति वर्ष है,
$100 \times 10^6$~\unit{m^2} पर: $10^8$~\unit{m^3}, यानी प्रति वर्ष \qty{0.1}{km^3}
भीतर आना चाहिए। (2) अंतर्वाह घटकर \qty{0.05}{km^3} रह जाता है: झील
$\qty{0.05}{km^3} = 5 \times 10^7$~\unit{m^3} खोती है, जो
$5 \times 10^7/10^8 = \qty{0.5}{m}$ की परत है। (3) तल गिरने के साथ किनारे
जल से बाहर निकल आते हैं और झील का क्षेत्रफल सिकुड़ता है: कम सतह, कम
वाष्पन, जब तक वाष्पन और अंतर्वाह फिर से संतुलित न हो जाएँ।
\end{solution}

\begin{solution}{exo:g10:freshwater-resources:14}
जलभृत $100 \times 0.20 = \qty{20}{m}$ जल रखता है, यानी
\qty{20000}{mm}: \qty{1}{mm} प्रति वर्ष पर \num{20000} वर्ष। सहारा
जलभृत का गहरा जल \num{200000} से एक मिलियन वर्ष पुराना है, दस से पचास गुना
अधिक: उसका पुनर्भरण अतीत की अधिक आर्द्र जलवायु में हुआ था।
\end{solution}

\begin{solution}{exo:g10:freshwater-resources:15}
$337/453\,000 \approx \qty{0.00074}{km}$, यानी लगभग \qty{0.74}{m} जल।
जल-तल \qty{4.8}{m} गिरा, लगभग साढ़े छह गुना अधिक। जल
चट्टान के केवल छिद्रों को भरता है (और उसका सारा जल बाहर नहीं निकलता): जब
भौम जलस्तर एक मीटर गिरता है, तो मुक्त हुआ जल एक मीटर से कहीं पतली परत होती है,
यहाँ लगभग $0.74/4.8 \approx \num{0.15}$~m।
\end{solution}

\begin{solution}{pb:g10:freshwater-resources:1}
\textbf{1.} $1093/68\,000 \approx \qty{0.016}{km}$, यानी लगभग \qty{16}{m}।

\textbf{2.} $1093 \times 10^9 \approx \qty{1.09e12}{m^3}$।

\textbf{3.} \qty{10}{g/L} का अर्थ \qty{10}{kg/m^3} है: $1.093 \times 10^{12} \times 10
\approx \qty{1.1e13}{kg}$, यानी लगभग \num{11} अरब टन।

\textbf{4.} $1093/71.1 \approx 15$ वर्ष।

\textbf{5.} उससे कोई नदी नहीं निकलती: उसका जल केवल वाष्प के रूप में (या
ज़मीन में रिसकर) निकल सकता है, जबकि लवण बना रहता है।

\textbf{6.} $71.1/68\,000 \approx \qty{0.00105}{km}$, यानी लगभग \qty{1.05}{m}।

\textbf{7.} $9.4/68\,000 \approx \qty{0.00014}{km}$, यानी लगभग \qty{0.14}{m}।

\textbf{8.} प्रति वर्ष $71.1 - 9.4 = \qty{61.7}{km^3}$।

\textbf{9.} $79 + 37 = \qty{116}{km^3}$: $61.7/116 \approx \qty{53}{\percent}$,
लगभग आधा।

\textbf{10.} आगत $1.0 + 6.2 = \qty{7.2}{km^3}$; निर्गत \qty{36.8}{km^3};
शुद्ध परिवर्तन $7.2 - 36.8 = \qty{-29.6}{km^3}$।

\textbf{11.} प्रति वर्ष $(502.7 - 327.2)/5 = \qty{35.1}{km^3}$: 1987 के बजट के
\qty{29.6}{km^3} के समान परिमाण-क्रम का (अलग-अलग वर्ष भिन्न होते हैं: 1984 में वाष्पन
\qty{47.9}{km^3} था)।

\textbf{12.} $36.8/0.00105 \approx \qty{35000}{km^2}$।

\textbf{13.} $35\,000/68\,000 \approx \qty{52}{\percent}$, लगभग आधा।

\textbf{14.} झील सिकुड़ चुकी थी: आधी सतह से लगभग आधा जल
वाष्पित होता है।

\textbf{15.} $345.6/1093 \approx \qty{32}{\percent}$, लगभग एक-तिहाई।

\textbf{16.} $10 \times 1093/345.6 \approx \qty{31.6}{g/L}$।

\textbf{17.} मापी गई लवणता कम है: लवण का एक भाग जल से निकल चुका था,
खुले तल पर क्रिस्टलित होकर पवन द्वारा लवणीय धूल के रूप में उड़ा ले जाया गया था।

\textbf{18.} प्रति वर्ष लगभग \qty{1.05}{m}, यानी वाष्पित मोटाई।

\textbf{19.} $7.2/35\,000 \approx \qty{0.0002}{km}$, यानी लगभग \qty{0.20}{m}।

\textbf{20.} $1.05 - 0.20 \approx 0.85$ (अपूर्णांकित मानों के साथ \num{0.84}):
\textbf{1987 में अरल सागर का तल लगभग \qty{0.8}{m} गिरा, एक ही वर्ष में
लगभग एक मीटर}।
\end{solution}
